\subsection{Double centering and a transport pivot over the integers}
\label{sec:integer-projections}

The following five equations use unrestricted integer variables and the
same six fixed coefficients throughout.  Each is a sum of squares of the
residuals in its accompanying JSON packet.
\begin{center}
\begin{tabular}{lrrr}
\hline
Packet & Variables & Residuals & Exact degree\\
\hline
\texttt{centered23\_degree10} &23&11&10\\
\texttt{pivot21\_degree16} &21&9&16\\
\texttt{pivot20\_degree20} &20&8&20\\
\texttt{pivot19\_degree32} &19&7&32\\
\texttt{pivot18\_degree88} &18&6&88\\
\hline
\end{tabular}
\end{center}
The inspected incoming Joint Guards package had the corresponding points
$23/12$, $21/20$, $20/32$, and $19/88$. The incoming Radical Guard Descent
$22/14$ form remains another useful tradeoff. These improvements do not
lower the ten-integer high-degree variable bound.

All new coordinate transformations below are integer-polynomial bijections
of the indicated parent zero sets. Consequently they preserve the parents'
ordinary, ZFC-formalizable equivalence with finite syntactic ZFC refutability.
Their unsolvability is therefore ZFC-provably equivalent to
$\operatorname{Con}(\mathrm{ZFC})$. This is an ordinary mathematical
composition; no new end-to-end kernel certificate or object-language ZFC
derivation is claimed.

The constituent two-root centering appears in incoming Budget and Centering,
and the monic transport pivot appears in incoming Round11. The new result is
their application to the newer Joint Guards systems. No historical priority
or global optimality claim is made.

\subsection{Fixed coefficients}
The coherent tuple is
\[
 \beta=2^{1064921}-1,\quad K_0,M_C,M_F,\quad
 A_*=2^{1113},\quad B_*=2^{666104}.
\]
Here $K_0,M_C,M_F$ are exactly the source-34
\texttt{mixed-parity-unrounded-1064921} native tuple, with $M_F$ already
shifted once. The six actual minimal unsigned little-endian byte strings
are supplied, compressed with gzip, with a complete binding manifest. Their
combined raw size is $615{,}883$ bytes. These are fixed numeral leaves,
not quantified parameters.

\subsection{The 23-variable equation}
Take independent integer coordinates
\[
 (J,F,\alpha,z,f,h,r,v,w,\tau,\eta,\zeta,\nu,\gamma,o,j,E,a,I,t,p_1,p_2,p_3).
\]
Define polynomial abbreviations only:
\begin{align*}
 q&=\beta(3J+1)+1,& W&=4(3w+1),& X&=qW,\\
 Y&=a-E,& k&=\eta+\zeta,& c&=Yk+\eta,\\
 V&=fo-c,& C&=q-F-\alpha,& H&=4a+3,\\
 \Delta&=a^2+H,& T&=E(c-\eta),& S&=X+\gamma H.
\end{align*}
Also put
\begin{align*}
 R&=((q-1)(q-F)+\alpha)(q^2-1)+(M_C+qM_F)(3J+1),\\
 G&=(A_*F)(B_*C)-(q-A_*F-B_*C)^2,\\
 Z&=\eta\zeta,\qquad
 \mathrm{Pos}=1+p_1^2+p_2^2+p_3(p_3+1).
\end{align*}
The eleven residuals are
\begin{align*}
 B_1&=a-E-(2v+1)q^3,\\
 B_2&=E-XY,\\
 B_3&=\tau^2+(2\tau-k)T-1,\\
 B_4&=2acS+S^2-Hc^2-1,\\
 B_5&=k-1-hE-R,\\
 B_6&=(K_0+W)C-F-z(q-1),\\
 B_7&=I-(2r+1)c^2,\\
 B_8&=I^2-\Delta(f^2-1),\\
 B_9&=(V+\nu)^2-I^2\nu(2V+\nu)-1,\\
 B_{10}&=V+k-1-cj-hE,\\
 B_{11}&=GZ-(t^2-G-Z+1)^2-\mathrm{Pos}.
\end{align*}
Thus the complete ordinary polynomial equation is
\[
 \boxed{\displaystyle P_{23}=\sum_{i=1}^{11}B_i^2=0.}
\]
Every displayed abbreviation is inlined in the supplied packet; it
introduces neither a variable nor a side condition.

For the proof, the incoming 25-variable octic has independent $c,V,y$ and
uses $k_0=1-V+jc+hE$ in its $c$ equation and first norm. It retains the
row $L=k-k_0$. Replace $k_0$ by $k$ using polynomial multiples of this
retained row. Then make the triangular root translations
\[
 \tau_{\mathrm{old}}=\tau+E(c-\eta),\qquad
 y_{\mathrm{old}}=V+\nu,
\]
and eliminate the two monic definitions $c=(a-E)k+\eta$ and $V=fo-c$.
The three cancellation identities are
\begin{align*}
 (\tau+T)^2-T(k+T)-1
   &=\tau^2+(2\tau-k)T-1,\\
 (ac+S)^2-(a^2+H)c^2-1
   &=2acS+S^2-Hc^2-1,\\
 I^2\bigl(V^2-(V+\nu)^2\bigr)+(V+\nu)^2-1
   &=(V+\nu)^2-I^2\nu(2V+\nu)-1.
\end{align*}
In the literal parent order, its thirteen restored rows are exactly
\[
 B_1,B_2,(a-E)B_{10},B_3+TB_{10},B_4,B_5-B_{10},
 B_6,0,B_7,B_8,B_9,B_{10},B_{11}.
\]
This gives both directions of the zero-set equivalence without any
division. The inverse root translations and coordinate projection are
unique. The residual degrees are
\[
 4,3,4,5,4,2,5,4,5,3,4.
\]
Only $B_7$ contains $r$, and
\[
 [r^2a^4\eta^4]P_{23}=4.
\]
Hence the total degree is exactly ten on the actual fixed-coefficient slice.

\subsection{A transport pivot that costs no dominant degree}
All four remaining parents contain
\[
 (K_0+W)(q-F-\alpha)-F-z(q-1)=0.
\]
Change the complement coordinate by $C=q-F-\alpha$, then solve the
now monic row for $F$:
\[
 F=(K_0+W)C-z(q-1),\qquad \alpha=q-F-C.
\]
Delete that row and $F$, and substitute everywhere else. Both the map and
its inverse use only integer polynomials, so this is a bijection of integer
zero sets. The packing register becomes
\[
 R=(q-F)(q^3-q)-C(q^2-1)+(M_C+qM_F)(3J+1),
\]
of degree at most five. The determinant $G$ has degree four, and its
contribution to the joint-sign residual has degree at most eight.
The larger Pell/sign residuals therefore already pay the resulting degree
cost in the selected parents. This yields the $21/16$, $20/20$, $19/32$,
and $18/88$ packets in the table.

\subsection{The complete 18-variable endpoint}
Take independent integer coordinates
\[
 (J,C,z,f,h,r,v,w,\tau,y,\gamma,o,j,c,t,p_1,p_2,p_3).
\]
In addition to $q,W,X$ and the transport pivot above, define
\begin{align*}
 Y&=(2v+1)q^3,& E&=XY,& a&=(X+1)Y,\\
 V&=fo-c,& k&=1-V+jc+hE,& \eta&=c-Yk,\\
 \zeta&=k-\eta,& H&=4a+3,& \Delta&=a^2+H,\\
 I&=(2r+1)c^2,& T&=E(c-\eta),& D&=X+ac+\gamma H.
\end{align*}
Use the same $R,G,Z,\mathrm{Pos}$ definitions. The six residuals are
\begin{align*}
 A_1&=\tau^2-T(k+T)-1,\\
 A_2&=D^2-\Delta c^2-1,\\
 A_3&=cj-V-R,\\
 A_4&=I^2-\Delta(f^2-1),\\
 A_5&=I^2(V^2-y^2)+y^2-1,\\
 A_6&=GZ-(t^2-G-Z+1)^2-\mathrm{Pos}.
\end{align*}
The equation is
\[
 \boxed{\displaystyle P_{18}=\sum_{i=1}^{6}A_i^2=0.}
\]
These six residuals are completely inlined in its packet, so the count
18 includes all quantified variables. Retaining $\eta$ as a nineteenth
independent coordinate and adding $c-\eta-Yk=0$ gives the degree-32 sibling.

\subsection{Degree certificates and verification}
The four pivot equations have the following exact nonzero coefficients:
\begin{align*}
 [a^4f^8o^4]P_{21}&=1,\\
 [f^4v^4J^{12}]P_{20}&=2^4 3^{12}\beta^{12},\\
 [\eta^4h^4w^4v^4J^{16}]P_{19}&=2^{12}3^{20}\beta^{16},\\
 [h^8w^8v^{16}J^{56}]P_{18}&=2^{32}3^{64}\beta^{56}.
\end{align*}
For the last two, the leading forms are
\[
 E_{\mathrm{top}}=1944\beta^4J^4wv,
 \qquad \eta_{\mathrm{top}}\big|_{P_{18}}
   =-104976\beta^7hJ^7wv^2.
\]
In $P_{19}$, the dominant joint-sign part is
$-(\eta hE_{\mathrm{top}})^2$ before the outer SOS square. In $P_{18}$,
it is $-\eta_{\mathrm{top}}^4$. These are the unique contributions at
the respective highest residual degrees. The first two certificates come
from unchanged reduced-auxiliary and strong-Pell rows; the pivoted guard
cannot contribute their selected monomials. Full derivations are given in
the companion proof.

The generator authenticates declaration closure, exact degrees and all
leading coefficients. A separate checker imports no builder: it checks
all thirteen centering identities, every literal pivot identity, the actual
six coefficient byte strings and additional signed modular transcription
diagnostics. The finite tests certify only the algebra and the data stated;
the universal native/Pell/refutation calibration remains an explicitly
inherited ordinary theorem.

The next concrete targets are a positivity certificate that avoids squaring
the degree-eight sign register in $P_{19}$, and an orientation theorem that
avoids the degree-eleven $\eta$ substitution in $P_{18}$. Either would
improve these equations more than further cosmetic factorization.
